\documentclass[conference]{IEEEtran}
\usepackage[T1]{fontenc}
\usepackage{amsmath,amssymb}
\usepackage{booktabs}
\usepackage{multirow}
\usepackage{graphicx}
\usepackage{xcolor}
\usepackage[hidelinks]{hyperref}
\usepackage{url}

% Open issues for the authors, and numbers still to be filled from the runs.
\newcommand{\kunal}[1]{\textcolor{blue}{[Kunal: #1]}}
\newcommand{\res}[1]{\textcolor{red}{#1}}
\newcommand{\sys}{HASHIRU}
\newcommand{\bench}{PURSE}

\begin{document}

\title{Agents That Hold the Money:\\Measuring How LLM Agents Spend What They Are Given}

\author{\IEEEauthorblockN{Kunal Pai, Parth Shah, Harshil Patel}
\IEEEauthorblockA{\kunal{affiliations and emails}}}

\maketitle

\begin{abstract}
Agent software now lets a language model fund the agents it spawns: a parent hands a child a slice of its budget,
and the child may hire in turn. Whether such agents spend well has not been measured. The budget-aware systems that
are evaluated split the money with an optimizer outside the agents, and the runtimes that let a model fund its
sub-agents do not evaluate the split. We present two artifacts. \sys{} v2 is a multi-agent system in which every agent
holds a pool reserved from its hirer's (dollars, plus measured memory for local models) and spends it by hiring cloud
models, local models and priced remote agents over open protocols (A2A for agents, MCP for tools); firing an agent
returns what it did not spend, and the ledger can be rebuilt from the trace alone. \bench{} is a benchmark built so
that the best split of the money is known: projects of eight jobs with exact answers, a dial for how much of each job's
difficulty is hidden until it is tried, outside agents to buy from, and calibrated outcomes for every job on every
option. On \bench{}, \res{[headline result: agent allocation against the knapsack and the up-front LLM split at equal
budgets, with the oracle gap]}. \sys{} v2 and \bench{} are source-available.
\end{abstract}

\section{Introduction}
\label{sec:intro}

Teams of language-model agents cost money, and how a team's budget is split decides what the team can do. A
hierarchical system (a CEO agent that hires managers and workers, as in \sys{} v1~\cite{pai2025hashiru}, or the
orchestrators of AutoGen~\cite{wu2024autogen}, MetaGPT~\cite{hong2024metagpt} and Magentic-One~\cite{fourney2024magenticone})
spends on every hire: a strong cloud model for a hard step, a cheap one for an easy step, an outside agent that sells an
exact answer. A team that spends its budget on the wrong steps fails on the ones that needed it, and a team that cannot
bound its spending cannot be trusted with a budget at all.

Agent software has started to hand this money to the agents. Automaton's \texttt{fund\_child} transfers credits to a
child agent~\cite{conway2026automaton}; wayland-core's \texttt{delegate} gives a sub-agent a dollar ceiling charged to
it and its descendants~\cite{ferrox2026waylandcore}; AgentLance lets a manager model set budgets for subtasks within
its payment~\cite{liu2026markets}. Budget-aware systems that are evaluated split the money with an optimizer outside the
agents instead (BAMAS~\cite{yang2026bamas}, ZEBRA~\cite{hamri2026zebra}), and benchmarks of budget awareness
(BAGEN~\cite{lin2026bagen}, $R^3$-Bench~\cite{wang2026r3bench}, EcoAgent-Bench~\cite{wu2026ecoagent}) measure one agent
spending its own tokens or actions. A ceiling that is checked after a call crosses it does not keep the call from
crossing it, and none of these systems is measured against the best split its money could buy. No arrangement covers
the case where agents hold the money for a whole hiring tree, every call is bounded by what its agent holds before it is
made, and the resulting allocation is measured against a known optimum.

To address these limitations, we built \sys{} v2. Every agent is its own A2A service, every tool is served over MCP,
and every hire reserves money from its hirer's pool. The hub, not the prompt, keeps spending safe: a call's output is
capped to what its agent has left, a call that cannot pay for an answer is refused with its price, firing returns what
was not spent, and an agent with one step left answers with what it has. Because every reservation, call and refund is
in the trace, the allocation can be measured, and we measured it on a benchmark built so that the best split of the
money is known.

\begin{keybox}{What this whitepaper contributes}
\begin{itemize}
  \item \textbf{\sys{} v2, an economy for agents:} each agent holds a pool reserved from its hirer's, in dollars and,
  for local models, measured memory; it declares how hard a job is and the hub hires the cheapest model that fits, or
  it buys from a priced remote agent or a whole remote team over A2A (\S\ref{sec:system}).
  \item \textbf{Budget safety by construction:} reserved pools with refunds, calls bounded by their pool before they
  are made, and a closing turn that turns an empty pool into an answer rather than a failure; on live projects, none
  of 48 hires spent past its reservation and every run answered (\S\ref{sec:pools}, \S\ref{sec:evaluation}).
  \item \textbf{Agents measured against optimizers:} on projects whose best split is known, \sys{}'s agents scored
  above sending every job to the strongest model at 37\% of its spend, matched a team that assigns tiers by level at
  35\% of its spend, and bought, escalated and re-reserved as the market priced their work (\S\ref{sec:evaluation}).
\end{itemize}
\end{keybox}

Each of the three safety mechanisms earns its place in the measurements. Run as first built, without the call bound
and the closing turn, the same agents overran a reservation by up to 372 times and lost half their runs to an empty
pool; with both, they lost none and spent 71\% of their budget. The remaining distance to the best split, and to a solver that
knows each job type's success rate in advance, is in how the agents judge difficulty, which is where \sys{} goes next
(\S\ref{sec:conclusion}).

\section{Who Splits the Money Today}
\label{sec:related}

We group earlier work by who decides how the money is split, and when (Table~\ref{tab:related}).

\textbf{An optimizer outside the agents.} BAMAS~\cite{yang2026bamas} chooses the models with an integer program and how
the agents work together with reinforcement learning, before the task starts. MasRouter~\cite{yue2025masrouter} sends
each query through a learned controller. ATOM~\cite{zhao2026atom} predicts how hard a query is and sets a budget for
it, and MALBO~\cite{sabbatella2025malbo} and SC-MAS~\cite{zhao2026scmas} tune who is on the team.
ZEBRA~\cite{hamri2026zebra} splits dollars across the steps of a fixed pipeline with a solver before anything runs, and
its language-model baseline loses to a knapsack solver. Mixture-of-Models~\cite{pecerskis2026mixtureofmodels} runs
preset team profiles. In all of these, the agents doing the work do not hold the money. \bench{} includes this setup as
a baseline (a knapsack solver over measured success rates, \S\ref{sec:purse-allocators}).

\textbf{Routing and cascades.} Model routers pick one model per query. RouteLLM~\cite{ong2025routellm} learns from
human preference data, Hybrid LLM~\cite{ding2024hybrid} picks between a small and a large model by predicted
difficulty, and FrugalGPT~\cite{chen2024frugalgpt} tries models from cheap to expensive until an answer is accepted.
RouterBench~\cite{hu2024routerbench} scores routers on recorded outputs against the best possible router, the same
method \bench{}'s replay uses. A router prices one query at a time. \bench{} prices a project whose jobs share one
budget, so money spent on one job is money another job cannot use.

\textbf{Agents that fund agents.} Agent runtimes released in 2026 let a model fund the agents it starts. Automaton sends
credits to a child agent running the same runtime~\cite{conway2026automaton}, and wayland-core's \texttt{delegate}
tool gives a sub-agent a dollar limit out of what its parent has left, charged to the sub-agent and the agents below
it~\cite{ferrox2026waylandcore}. AgentLance~\cite{liu2026markets} lets a manager model set enforced budgets for
subtasks, one level deep, with a market choosing the workers, and evaluates the whole system on benchmarks. \sys{}
v1~\cite{pai2025hashiru} came before these, with a CEO that hires and fires agents under cost and memory budgets, one
level deep. None of them measures how well the agents divide the money.

\textbf{Budget awareness in one agent.} BATS~\cite{liu2026budgetaware} tells a tool-using agent how much budget it has
left. BAGEN~\cite{lin2026bagen} finds that strong agents are not always budget-aware, and are too optimistic about tasks
they cannot finish with the money left. $R^3$-Bench~\cite{wang2026r3bench} has one model split its own output tokens or
actions across six problems of hidden difficulty, and scores it against the best possible split, computed offline from
how each problem responds to more effort. \bench{} computes its best split the same way. EcoAgent-Bench~\cite{wu2026ecoagent}
prices an agent's actions (lookups, searches, a stronger model, a human), AllocBench~\cite{wang2026allocbench} measures
investment in reusable tools, and CostBench~\cite{liu2026costbench} measures planning for the lowest cost and
replanning when tools fail or prices change. In all of them, one agent spends on itself.

\textbf{Local and cloud models, and agent economies.} Minions~\cite{narayan2025costefficient} and
PAPILLON~\cite{li2025papillon} pair a local model with a cloud model under a fixed protocol. Virtual agent
economies~\cite{tomasev2025virtual} and Magentic Marketplace~\cite{bansal2025magentic} study markets between agents.
Payment protocols (AP2~\cite{google2026agent}, the A2A x402 extension~\cite{googleagenticcommerce2025a2a}) move the
money but do not decide how to split it. \sys{} v2 uses the open protocols these build on, MCP for
tools~\cite{modelcontextprotocol2026specification} and A2A for agents~\cite{a2aproject2026agent2agent}.

\begin{table}[!htb]
\centering
\caption{Who splits the money, and when. Like the 2026 agent runtimes, \sys{} v2 lets agents hold budgets. Unlike
them, it caps each call before it is made and is compared with the best possible split.}
\label{tab:related}
\footnotesize
\begin{tabular}{@{}p{4.2cm}p{3.6cm}p{2.6cm}p{3.6cm}@{}}
\toprule
\tablehead Line of work & Who splits the money & When & Compared with the best split \\
\midrule
BAMAS, ZEBRA, MasRouter & optimizer & before the task & no \\
RouteLLM, FrugalGPT & router, cascade & per query & best router \\
Automaton, wayland-core & the model & during & no \\
AgentLance & manager model & per subtask & no \\
$R^3$-Bench, EcoAgent & one agent, on itself & during & yes \\
\textbf{\sys{} v2} & \textbf{the agents, to depth 2} & \textbf{during, each call capped} & \textbf{yes} \\
\bottomrule
\end{tabular}
\end{table}

\sys{} v2 shares agent-held budgets with these runtimes, and borrows $R^3$-Bench's way of computing the best split, but
it differs in two ways. (1) \emph{Limits before the call.} A hire sets money (and memory) aside that firing refunds, and
each call's output is capped at what its agent has left before the call is made. wayland-core's sub-agent, by
contrast, stops itself only after it has gone over its limit. (2) \emph{Compared with the best split.} To our
knowledge, it is the first system whose agent-held budgets, across a hiring tree, are compared with a solver, an
up-front plan and fixed teams at the same budget.

\section{\sys{} v2: An Economy for Agents}
\label{sec:system}

\sys{} v2 rebuilds v1~\cite{pai2025hashiru} on open protocols. A user's request goes to a CEO agent, which may hire
\emph{managers} (agents that use tools and may hire) and \emph{workers} (agents that answer once), to a configurable
depth (two levels below the CEO by default). Every agent is its own A2A service with its own Agent Card, every tool is
served over MCP, and \sys{} itself is reachable as an OpenAI-compatible model, an A2A agent and an MCP server. One
process serves everything on one port. Every model call, tool call and message is written to a JSONL trace.

\subsection{From v1 to v2}
v1 had one level of hierarchy, a hand-written table of resource units, and no analysis of the local-against-cloud
trade (its own stated limitations). v2 lets managers hire with money they were given, prices every cloud call from
LiteLLM's price map~\cite{berriai2026litellm} (local calls cost no dollars but hold the memory their model really uses,
measured from Ollama~\cite{ollama2026ollama}), and is measured by \bench{}. Rules that v1 left to prompts (who may
spend what, where private data may go) are enforced in code: in our experience with an earlier prototype, rules stated
only in prompts were broken.

\subsection{Pools, reservations and refunds}
\label{sec:pools}
Each agent $a$ holds a reservation $R_a$ carved from its hirer's pool. What it can still spend is
\begin{equation}
\mathrm{left}(a) = R_a - \mathrm{spent}(a) - \sum_{c \in \mathrm{hires}(a)} \mathrm{hold}(c),
\label{eq:left}
\end{equation}
where a live hire holds its whole reservation and a fired one holds only what it (and everyone below it) actually
spent. Firing therefore returns the unspent part, and two hires running at once cannot spend the same dollars: each
holds its reservation from the moment it is hired. Table~\ref{tab:pool} works one session through.

\begin{table}[t]
\centering
\caption{One session's pools (dollars). The CEO starts with \$0.50.}
\label{tab:pool}
\footnotesize
\begin{tabular}{@{}lrr@{}}
\toprule
Step & CEO left & m1 left \\
\midrule
Session opens, \texttt{budget\_usd} = 0.50 & 0.50 & \\
CEO spends 0.01 thinking & 0.49 & \\
CEO hires manager m1 (0.20) and worker w1 (0.10) & 0.19 & 0.20 \\
m1 hires w2 (0.05) and spends 0.02 & 0.19 & 0.13 \\
w1 spends 0.03, then the CEO fires it & 0.26 & 0.13 \\
\bottomrule
\end{tabular}
\end{table}

The checks are in code, not in prompts: a hire cannot reserve more than its hirer has left, an agent whose pool is
spent makes no further cloud call, and a remote agent is paid from its hirer's reservation, which must cover its
declared price before each message. A local hire is refused if its model does not fit in free memory (or in a cap,
\texttt{memory\_mb}). Separately, \texttt{max\_usd} caps real spend for the whole process. Because every reservation,
call and fire is in the trace, the ledger can be rebuilt from the trace alone, so budget rules can be checked after the
fact.

\subsection{Hiring by declared difficulty}
A hiring agent may name a model, or declare how hard a job is (easy, medium or hard) and let the hub choose. The hub
turns a difficulty into a list of models, best first, and hires the first that the modes, the hirer's pool and the
data's sensitivity allow: an easy job tries small local models, then cheap cloud models; a hard job the dearest cloud
models first. A worker that fails or answers empty moves up its list and gets the message again, up to a limit
(\texttt{max\_escalations}, two by default). Which models exist, and their tiers, come from the machine itself:
Ollama's pulled models that call tools, and every cloud model in LiteLLM's map whose provider key is set, tiered by
price or by size. Declared difficulty is the agent's own judgment, so \bench{} measures it against the hidden
difficulty of each job (\S\ref{sec:evaluation}).

\subsection{Make or buy over open protocols}
Besides models, a manager may hire a remote A2A agent as a worker (\texttt{a2a/<name>}). The hub reads the agent's
card, prices it through a small A2A extension (a card declares a price per message or per token; a reply reports what
the request cost), charges the larger of the declared price and the reported cost to the hire's pool, and pins the
agent to its configured URL. Since a \sys{} hub is itself an A2A agent, one hub can hire another: a whole team, with
its own CEO and pools, appears to the hirer as one priced worker. This turns hiring into make or buy: an agent can make
an answer with a model it pays per token, or buy it from an outside agent at a price.

\subsection{Quotes, authority and sensitivity}
Three further mechanisms read naturally in economic terms. \emph{Plan mode is a quote:} the CEO may only read and
propose; the hub previews each proposed hire as it would be made (the model the router picks, what it reserves, the
memory it takes) and checks the totals against the pool, and approving the plan approves its hires. An approved plan's
steps compile to a graph (LangGraph~\cite{langchain2026langgraph}) that hires, runs each step once its inputs exist,
fires each hire after its last step, and lets the CEO answer once. \emph{Approvals are spending authority:} a hire or a
turn above a threshold waits for the user. \emph{Sensitivity limits where work may go:} private data never reaches a
cloud or remote model as it is; it is masked with placeholders on the way out and restored on the way back (Presidio's
recognizers~\cite{presidio2026presidio} and a local model find the values), or refused.

\subsection{Implementation}
\sys{} v2 is 6,323 lines of Python 3.13 with exact pins for its protocol and model libraries, and 446 tests that run
without API keys or Ollama. The two changes \bench{} needed (a switch that takes the web tools away, so that no answer
can be looked up, and the memory cap) were added test-first. Running \bench{}'s first real calls found one defect:
a reasoning model (\texttt{gpt-6-luna}) rejects temperature 0, so \sys{} now lets LiteLLM drop that parameter for such
models only.

\section{How We Measured}
\label{sec:purse}

To know whether agents spend well, we need tasks where the best split of the money is known. Our benchmark, \bench{},
builds such tasks. Each project is a set of jobs that share one budget. Every job has an exact answer, and every job has
been tried on every option in advance, so the best split for any budget can be computed. The score then depends on how
the money is split, and the gap to the best split shows how much a strategy left on the table.

\subsection{Jobs and projects}
A project has eight jobs, two from each of four job types (Table~\ref{tab:families}), and one budget $B$. Each job has
a visible level from 1 to 5 (its stated size), a hidden difficulty (low, mid or high), a value of 1, 2 or 4 points
chosen independently of difficulty, a standard prompt and an exact checker. A strategy reads every job with its level
and value and answers with one JSON object. The score is the share of points answered correctly. Jobs are generated
from random seeds, so they cannot be looked up, and \sys{}'s web tools are turned off in every run.

\begin{table}[!htb]
\centering
\caption{\bench{}'s four job types.}
\label{tab:families}
\footnotesize
\begin{tabular}{@{}p{2.2cm}p{4.0cm}p{5.2cm}p{3.0cm}@{}}
\toprule
\tablehead Job type & Set by the level & Set by the hidden difficulty & Can be bought from \\
\midrule
Arithmetic & operators (2 to 22), digits (2 to 8) & multiplications, digit sizes (carries) & \texttt{calc} (levels 1 to 4) \\
Logic grids & houses, attributes & guesses a constraint solver must make (0, 1, 2+) & \texttt{csp}; a free checker \\
Aggregation & records (40 to 2{,}500) & fields in the condition (1, 2, or 2 plus near misses) & no one \\
String pipelines & operations (2 to 10), length & alphabet, operations that depend on position & no one \\
\bottomrule
\end{tabular}
\end{table}

For the second round (\S\ref{sec:evaluation}), the hidden difficulties were spread further apart, so that knowing them
would be worth more. In arithmetic, the hidden difficulty also changes the numbers' digits (one fewer, the same, or two
more) and the share of multiplications (10\%, 35\% or 70\%). In aggregation, it sets how many records match (a handful,
about 1 in 20, or about 1 in 7) and triples the near misses. In strings, it halves, keeps or doubles the starting
string. At high difficulty all but one operation depends on position, and at low difficulty the Caesar shift is
dropped.

\subsection{How much of the difficulty is visible}
With probability $\rho$, a job's hidden difficulty is the one its visible level suggests (levels 1 and 2 low, 3 mid, 4
and 5 high). Otherwise it is random. At $\rho = 1$ the level tells the whole story, so a plan made up front with the
right statistics loses nothing to a strategy that adapts. At $\rho = 0$, difficulty can only be learned by trying. The
value of knowing the hidden difficulty,
\[
\mathrm{VOI} = \mathrm{score}(\text{oracle knowing the difficulty}) - \mathrm{score}(\text{oracle knowing only the levels}),
\]
is the most that any adaptive strategy can gain over the best up-front plan, and we report it for every setting. Each
strategy gets the statistics it would measure on projects with its own $\rho$, which are each difficulty's success
rates weighted by $P(\text{difficulty} \mid \text{level}) = \rho\,[\text{difficulty matches the level}] + (1 - \rho)/3$.

\subsection{The market}
\label{sec:purse-market}
The options are models at different prices (a cheap and an expensive reasoning model in the first round, plus a third,
\texttt{gpt-5-nano}, in the second) and outside agents reached over A2A. Two outside agents are exact programs.
\texttt{calc} does exact arithmetic (it is \sys{}'s example agent, and reads expressions of up to 200 characters, so it
refuses level 5), and \texttt{csp} solves logic grids. The second round adds \texttt{gridfast}, which sells the same
grids at a quarter of \texttt{csp}'s price and gets 40\% of them wrong (it swaps one attribute between two houses), so
that the cheapest seller is not always a good one. Outside agents charge a price per message, tested at
$c_\mathrm{ref} \times \{0.25, 1, 4, 16\}$, where $c_\mathrm{ref}$ is the cheap model's median cost on the jobs they can
do. A free \texttt{grid\_check} tool says whether a grid solution meets every clue, up to three times per puzzle
(checking is cheaper than solving, as with unit tests).

\subsection{Recorded results, replay and the oracle}
Before any strategy runs, every job is sent with its standard prompt to every option that can do it (once in the first
round, three times in the second), through a \sys{} hub (one hire, one message, one fire). We call this
\emph{calibration}. The \emph{results table} records each answer, whether it is right, and what it cost. A replayed
call returns one of the job's recorded attempts, and each seed starts from a different one. The table has three uses.
First, it gives the informed strategies their statistics, taken from the other half of the jobs, so a project never
sees statistics that include its own jobs. Second, it defines the \emph{oracle}. Oracle-full picks, for each job, the
option that gets the most points under the budget, using each job's recorded success and cost (a multiple-choice
knapsack, as in $R^3$-Bench~\cite{wang2026r3bench}). No strategy can beat it. Third, it lets any strategy that sends
standard prompts be scored by \emph{replay} on recorded results instead of new model calls, as RouterBench does for
routers~\cite{hu2024routerbench}. Each project's budget is $B = \tau\,C_\mathrm{full}$, where $C_\mathrm{full}$ is what
oracle-full spends at its best score when money is unlimited. So $\tau = 1$ is just enough money for the best answers,
and $\tau = 0.5$ is half of that.

\subsection{Strategies}
\label{sec:purse-allocators}
The main comparison crosses who splits the money (a language model or a solver) with when (up front, or during the
task). All four get the same project text, prices, statistics and budget (Table~\ref{tab:twobytwo}).

\begin{table}[!htb]
\centering
\caption{The four strategies in the main comparison.}
\label{tab:twobytwo}
\footnotesize
\begin{tabular}{@{}p{1.6cm}p{6.4cm}p{6.4cm}@{}}
\toprule
\tablehead  & Up front & During the task \\
\midrule
LLM & \textbf{LLM split.} The CEO plans one hire per job in plan mode, and the plan runs unchanged (like ZEBRA's LLM
baseline). & \textbf{\sys{}.} The CEO hires, asks, fires, rehires and moves up to better models as results come in. \\
Solver & \textbf{Knapsack.} Picks options to get the most expected points under the budget, using the statistics. &
\textbf{Cascade.} Takes the cheapest option expected to succeed, and moves up a tier when a check fails or two answers
disagree. \\
\bottomrule
\end{tabular}
\end{table}

Around them are oracle-full, fixed strategies (every job on one model; a hybrid that picks by visible level, in the
style of AgentCARD~\cite{jiang2026specialize}; an equal split; a random choice), and \sys{} variants (a CEO that names
models, as in v1, and a CEO that sees prices but no statistics).

\subsection{Validity checks}
\label{sec:checks}
Before any strategy runs, seven checks on the results table make sure \bench{} really measures how the money is split
(Table~\ref{tab:checks}). We tested them on made-up tables built to pass, and on one table per check built to fail
exactly that check. Check V8 (replay agrees with live runs) needs live runs.

Both tables behind \S\ref{sec:evaluation} fail three of the seven, and not the same three. On the first table, V1
failed narrowly (the oracle beat the best fixed strategy by 9.7 points, not 10). V4 failed because the expensive model
had room to improve in only one job type (long aggregation jobs), not two. V5 failed outright. Knowing each job's
hidden difficulty was worth 0.009 of the points at $\rho = 0$, against the 0.08 the check asks for. We built the second
table to fix these, with a third model, a cheap but unreliable grid seller, more spread-out difficulties and three
attempts per job. V1 now holds (10.7 points). V5 still fails, and by more (0.002). The wider spread did separate the
cheap models (\texttt{gpt-6-luna} answers 100\% of low-difficulty aggregation jobs and 43\% of mid ones). But knowing
the hidden difficulty is worth between $-0.007$ and 0.002 of the points at every budget we tried ($\tau$ from 0.25 to
2), and 0.01 to 0.02 with the expensive model removed. With only four jobs per group, each difficulty's rates rest on
twelve attempts, which are as noisy as the differences they are meant to show. V3 now fails because
\texttt{gpt-5-nano} is not simply a weaker model. It beats \texttt{gpt-6-luna} on arithmetic (53\% against 25\% at
level 5) and loses to it on aggregation and strings, so the three models cross over by job type instead of lining up,
and on logic grids the expensive model never leads the cheap one by 20 points. V4 fails as before. So neither table can
show what adapting during the task is worth, and the during-the-task column of Table~\ref{tab:twobytwo} cannot be
tested against the up-front one. Normally a failed check means a job generator is retuned and its jobs are measured
again. We report both tables as they are, and a table that passes V5 is on the roadmap (\S\ref{sec:conclusion}).

\begin{table}[!htb]
\centering
\caption{Validity checks on the results table, run before any strategy, and how the two tables behind
\S\ref{sec:evaluation} do.}
\label{tab:checks}
\footnotesize
\begin{tabular}{@{}p{3.0cm}p{6.7cm}p{2.4cm}p{2.4cm}@{}}
\toprule
\tablehead Check & Holds when & First table & Second table \\
\midrule
V1 Splitting matters & the oracle beats the best fixed strategy by $\geq$ 10 points at $\tau = 0.5$ & \textbf{fails} (9.7) & holds (10.7) \\
V2 The budget is tight & the expensive model alone covers $<$ 60\% of the points at $\tau = 0.5$ & holds (59.6\%) & holds (51.7\%) \\
V3 Models separate & for each job type, a level where the cheap model succeeds $\geq$ 80\% (and a small model, if present, $\leq$ 40\%), and one where the expensive model leads the cheap one by $\geq$ 20 points & holds (4 of 4 types) & \textbf{fails} (0 of 4) \\
V4 Room to improve & the expensive model succeeds $\leq$ 60\% at level 5 in $\geq$ 2 job types & \textbf{fails} (1 type) & \textbf{fails} (1 type) \\
V5 Hidden difficulty matters & VOI $\leq$ 0.02 at $\rho = 1$ and $\geq$ 0.08 at $\rho = 0$ & \textbf{fails} (0.000; 0.009) & \textbf{fails} (0.000; 0.002) \\
V6 Buying responds to price & the oracle's buying falls $\geq$ 40 points across the price sweep & holds (53 points) & holds (56 points) \\
V7 No shortcut & the most common answer per job type scores $\leq$ 10\% & holds (4.3\%) & holds (5.6\%) \\
\bottomrule
\end{tabular}
\end{table}

\FloatBarrier  % the validity checks table stays in this section

\section{Results}
\label{sec:evaluation}

We ask three questions. Do agents that hold the money stay within it? Do they use the market that \sys{} gives them?
And how well do they split the money?

\subsection{Setup}
We ran two rounds. The first (2026-11-01) built the guards and measured them. The second (2026-11-02) added a third
model and an unreliable seller, measured the models again on a wider set of jobs, and tested the agents on more
projects.

\paragraph{Models and prices} The CEO is \texttt{gpt-6-luna} (\$0.10 / \$0.50 per million input / output tokens). The
first round's market has \texttt{gpt-6-luna} (the cheap model), \texttt{muse-spark-1.3} from Meta's Model API (\$1.25 /
\$4.25, the expensive model), and the \texttt{calc} and \texttt{csp} agents. The second adds \texttt{gpt-5-nano}
(\$0.05 / \$0.40) and \texttt{gridfast}. Prices are LiteLLM's list prices on the run dates. \texttt{gpt-5-nano} is the
cheapest per token, but it reasons for longer, so one attempt on it costs three times as much as one on
\texttt{gpt-6-luna} (\$0.0030 against \$0.0010). We tried other small models and dropped them (Gemma 4 31B, whose
hosted endpoint did not answer, and Muse Glimmer 30B, which costs more per token than \texttt{gpt-6-luna}). We left out
local models, since on the laptop we had (Apple M5, 24 GB) a 4B thinking model took six minutes on one arithmetic job.

\paragraph{Calibration} The first round sent 120 jobs (two per job type, level and difficulty) once to each option, for
\$0.18 on \texttt{gpt-6-luna} and \$3.18 on \texttt{muse-spark-1.3}. The outside agents are our own programs and cost
nothing real. The second round sent 240 jobs (four per group) three times to each option, for \$0.70 on
\texttt{gpt-6-luna}, \$2.17 on \texttt{gpt-5-nano} and \$20.32 on \texttt{muse-spark-1.3}. Table~\ref{tab:calibration}
shows where the second round's models differ, and where they cross over.

\begin{table}[!htb]
\centering
\caption{Share of attempts answered right in the second round's calibration, by level (1 to 5; 36 attempts each).}
\label{tab:calibration}
\footnotesize
\setlength{\tabcolsep}{4pt}
\begin{tabular}{@{}lllll@{}}
\toprule
\tablehead Job type & \texttt{gpt-5-nano} & \texttt{gpt-6-luna} & \texttt{muse-spark-1.3} & Outside agent \\
\midrule
Arithmetic & 100 100 100 81 \textbf{53} & 94 100 97 \textbf{64 25} & 100 100 100 92 92 & 100 100 100 100 \textbf{0} (\texttt{calc}) \\
Logic grids & 94 100 100 100 94 & 100 94 94 94 94 & 100 100 100 100 100 & 100 at every level (\texttt{csp}) \\
 & & & & 67 67 50 58 50 (\texttt{gridfast}) \\
Aggregation & 100 92 \textbf{44 25 0} & 92 86 64 53 \textbf{36} & 100 100 92 100 \textbf{47} & none \\
String pipelines & 100 94 92 64 \textbf{47} & 92 83 94 86 67 & 100 100 100 94 100 & none \\
\bottomrule
\end{tabular}
\end{table}

\paragraph{What these tables can show} Three of the seven validity checks fail on each table (\S\ref{sec:checks}), and
each failure limits what this section can claim. V5 fails on both (knowing the hidden difficulty is worth at most 0.009
of the points at $\rho = 0$), so these runs cannot show what adapting during the task is worth, for any strategy. We
compare the up-front and during-the-task strategies of Table~\ref{tab:twobytwo} on what they score and spend, not on how
well they adapt. V4 fails, so the expensive model solves nearly every job it is given (it drops to 60\% or below only on
level-5 aggregation). V3 fails on the second table, so its three models do not line up from weak to strong, since
\texttt{gpt-5-nano} beats \texttt{gpt-6-luna} on arithmetic and loses to it on aggregation and strings.

\paragraph{Runs} The first round drew ten projects at $\rho = 0.5$ and $\tau = 0.5$ (mean budget \$0.056), and ran
\sys{} v2 live twice on each of the first six (mean budget \$0.043). Its outside agents charged \$0.001 a message live,
against the $c_\mathrm{ref}$ of \$0.00023 that the replayed strategies paid. The second round draws twenty projects with
the same $\rho$ and $\tau$ (mean budget \$0.058) and runs every strategy three times on each. \sys{} v2 and the LLM split
use three seeds, and the replayed strategies use three of each job's recorded attempts. Its outside agents charge
$c_\mathrm{ref}$ (\$0.000215 a message, and a quarter of that for \texttt{gridfast}), live and in replay alike. The
solvers and fixed strategies are replayed on the results table. The LLM split makes its plan live and carries it out by
replay. Each is scored under two budget rules. Under the first hub's rule, a call only needs some money left and is
charged in full, even past the budget. Under v2's rule, a call is capped at what is left before it is made, as it is for
\sys{} v2. Replay cannot know what a capped call would have answered, so it counts it as no answer. In the first
round's live runs, all ten calls by hires that hit their cap returned nothing.

We threw away two attempts at the second round's live runs. The first sent nine processes at once, went over the
provider's limit of two million tokens a minute, and lost runs to refused calls. We made a rate-limited call wait and
retry (for up to 104 seconds), and no call in the runs we report failed for this reason. In the second attempt, the
project text put the market, the statistics and the answer format after the last job, inside that job's section, which
\texttt{--by-reference} hides from the CEO. The text now puts them before the jobs, for every version. In total we spent
\$33.21. The first round cost \$4.50. The second cost \$23.42 to calibrate (\$20.32 of it on \texttt{muse-spark-1.3},
and \$0.23 on pilots), \$4.30 for its live runs and \$0.99 for the two discarded attempts.

\subsection{What we found}
Overall, our experiments show the following results:

\textbf{The guards cap every hire.} The second round made 120 live runs of \sys{} v2 (two versions, twenty projects,
three seeds) with 551 hires. No hire spent more than 1.11 times its budget. Thirteen went over, by 9 to 11\%, each on a
single call, and all thirteen were aggregation jobs, whose inputs are long lists of records. A call's output is capped at
what the budget has left after the estimated cost of its input, so an overrun can only happen when the hub undercounts
the input's tokens by more than its 10\% margin. In the first round, none of 48 hires went over (Table~\ref{tab:ablation}).

\textbf{Each guard is needed.} In the first round we ran the same projects with \sys{} as first built (without the call
cap and the last turn), and with the call cap alone (Table~\ref{tab:ablation}). Without the cap, a budget was a limit
checked before a call, and the call's cost was charged after it. Hires given \$0.0005 spent up to 372 times that, the
overruns were charged to the CEO's budget, and three of six runs ended with no answer after spending 153\% of the
budget. With the cap alone, no hire went over, but the CEO reread every job, long documents included, on every turn,
and four of six runs ended without an answer (three because the CEO could not pay for its own answer, one at its
40-turn limit). With the last turn added, two of twelve runs were lost and the score rose to 0.566.

\begin{table}[!htb]
\centering
\caption{\sys{} with and without its guards, on the first round's six projects.}
\label{tab:ablation}
\small
\begin{tabular}{@{}lrrr@{}}
\toprule
\tablehead & As first built & Capped calls only & \sys{} v2 \\
\midrule
Runs & 6 & 6 & 12 \\
Mean score & 0.374 & 0.278 & \textbf{0.566} \\
Hires that spent past their budget & 9 of 26 & 0 of 25 & \textbf{0 of 48} \\
Largest overrun (spent / given) & 372$\times$ & 1.0$\times$ & \textbf{1.0$\times$} \\
Runs that ended without an answer & 3 of 6 & 4 of 6 & \textbf{2 of 12} \\
Spent / budget & 1.53 & 0.62 & \textbf{0.71} \\
Calls refused with a price quote & 0 & 19 & 25 \\
\bottomrule
\end{tabular}
\end{table}

\textbf{Keeping long inputs away from the CEO raises the agents' score by 0.136.} The second round ran \sys{} v2 as it
was, and with two changes. Long sections were hidden from the CEO (\texttt{request\_by\_\allowbreak reference}), and hires were
chosen by the success rates the CEO is shown (\texttt{success\_rates}). With both, the agents scored 0.599 against
0.463, a difference of $+0.136$ over the twenty projects (95\% interval $[+0.067, +0.216]$). The CEO reads the request's
outline (8.7 thousand characters on the first project) instead of the whole request (264 thousand), and its own calls
fell from 43\% of the budget to 32\%. Messages to outside agents now carry the puzzle the CEO names, so \texttt{csp}
answered 48 of its 49 messages right, where before, 24 of its 44 replies had been ``no grid puzzle in the text''. The
outside agents together answered 74 of 119 right, against 30 of 122. Runs where the CEO could not afford to read its
request once fell from 12 to 0, and lost runs from 17 to 11. We did not test the two changes separately. Every effect
above comes from hiding long sections, and what choosing by success rates adds on its own is not measured.

\textbf{With enough budget to run the CEO, the agents match the cheap model alone.} Table~\ref{tab:main} scores every
strategy under v2's budget rule, with the first rule's scores beside them. On the sixteen projects whose budget is
\$0.006 or more, the agents scored 0.737. That is about the same as \texttt{gpt-6-luna} alone (0.740; difference
$-0.003$, interval $[-0.038, +0.030]$) and the equal split (0.761; $-0.024$, $[-0.056, +0.007]$), and better than
random choice ($+0.067$, $[+0.005, +0.134]$). The agents also scored above the cascade, the LLM split and
\texttt{gpt-5-nano} alone, but by margins too small for sixteen projects to settle. They trail the knapsack ($-0.080$,
$[-0.129, -0.023]$), and they spent 3.9 times what the cheap model alone spent.

\textbf{The agents win on hard jobs that a specialist can do.} Two job types have an outside agent in the market
(\texttt{calc} for arithmetic, \texttt{csp} and \texttt{gridfast} for logic grids). On the sixteen projects, over the
hard jobs of those two types (26 jobs, three runs each), the agents answered 62 of 78 right (79\%), against 54 of 78
for \texttt{gpt-6-luna} alone (69\%). That is 24 against 18 on hard arithmetic, and 38 against 36 on hard grids. They
were ahead on six projects and behind on two, and the interval over projects only just reaches zero. They send grids to
\texttt{csp}, which answered all 41 of its messages on these projects right, and they give hard jobs a second hire more
often than easy ones (46 of the 81 jobs that went to more than one hire were hard, while hard jobs are 37\% of all
jobs). On the other jobs they do worse than the cheap model (226 of 306 right, 74\%, against 239, 78\%), and the two
cancel out to the matching overall score above.

\textbf{Below a minimum budget, the agents cannot work.} Four of the twenty projects have budgets of \$0.0009 to
\$0.0051, because their jobs are cheap to solve and $B = \tau\,C_\mathrm{full}$. On them, \texttt{gpt-6-luna} alone
scores 0.724, and the LLM split, whose one planning call comes out of the budget, scores 0.635. The agents score 0.046.
A CEO that hires, asks and fires pays for a turn at each step. All eleven runs the agents lost were on these four
projects, where the CEO came back empty after spending most or all of the budget on its own calls. On the other sixteen
they lost none of 48. Over all twenty projects, the agents' 0.599 is below the cheap model alone (0.736).

\begin{table}[!htb]
\centering
\caption{Strategies on the second round's twenty projects ($\rho = 0.5$, $\tau = 0.5$, mean budget \$0.058), sorted by
score under v2's budget rule, and on the sixteen projects whose budget is \$0.006 or more. Each runs three times per
project (sixty runs). ``Vs.\ Luna'' is spending divided by that of every job on \texttt{gpt-6-luna} (\$0.0079 a
project). The last column replays the same strategies under the first hub's rule, which charges a call in full even
past the budget. \sys{} v2's two versions run live, as it was and with long sections hidden (both during the task).}
\label{tab:main}
\footnotesize
\setlength{\tabcolsep}{3.6pt}
\begin{tabular}{@{}lrrrrrrr@{}}
\toprule
\tablehead Strategy & Score & Sixteen & Spent (\$) & Of budget & Vs.\ Luna & Own calls (\$) & First rule \\
\midrule
Oracle-full (knows outcomes) & 0.826 & 0.838 & 0.028 & 0.48 & 3.50$\times$ & 0 & 0.830 (0.49) \\
Knapsack (solver, up front) & 0.788 & 0.817 & 0.035 & 0.60 & 4.42$\times$ & 0 & 0.794 (0.68) \\
Every job on \texttt{gpt-6-luna} & 0.736 & 0.740 & 0.008 & 0.14 & 1.00$\times$ & 0 & 0.723 (0.15) \\
Equal split & 0.713 & 0.761 & 0.011 & 0.18 & 1.35$\times$ & 0 & 0.713 (0.18) \\
LLM split (LLM, up front) & 0.695 & 0.709 & 0.033 & 0.57 & 4.15$\times$ & 0.006 & 0.699 (0.69) \\
Cascade (solver, during) & 0.682 & 0.712 & 0.050 & 0.86 & 6.29$\times$ & 0 & 0.762 (1.58) \\
Every job on \texttt{gpt-5-nano} & 0.628 & 0.700 & 0.023 & 0.39 & 2.87$\times$ & 0 & 0.630 (0.41) \\
Random & 0.612 & 0.670 & 0.029 & 0.50 & 3.69$\times$ & 0 & 0.614 (1.05) \\
\textbf{\sys{} v2, long sections hidden} & \textbf{0.599} & \textbf{0.737} & \textbf{0.031} & \textbf{0.52} & \textbf{3.84$\times$} & 0.006 & live \\
Hybrid by level & 0.493 & 0.494 & 0.052 & 0.90 & 6.60$\times$ & 0 & 0.531 (2.14) \\
\sys{} v2 as it was & 0.463 & 0.579 & 0.036 & 0.61 & 4.47$\times$ & 0.025 & live \\
Every job on \texttt{muse-spark-1.3} & 0.342 & 0.427 & 0.054 & 0.93 & 6.82$\times$ & 0 & 0.447 (2.04) \\
\bottomrule
\end{tabular}
\par\smallskip{\footnotesize\raggedright First rule: score, with spent / budget in parentheses. ``Sixteen'': score on the
sixteen projects whose budget is \$0.006 or more.\par}
\end{table}

\textbf{The agents respond to prices.} A hire's call that its budget cannot pay for is refused before it reaches the
provider, and the refusal quotes the price. In the second round, 203 calls were refused this way in the version that
hides long sections (253 in the other). In 112 of them, a later hire in the same run was given more than the refused one
(156 in the other). Table~\ref{tab:refusal} follows one such exchange from a first-round log. The CEO called a level-4
logic grid easy and gave it to \texttt{gpt-6-luna} with \$0.001. The hire's first call reasoned up to the cap its budget
allowed and returned nothing. Its next call was refused, and so was the same call on the expensive model after the hub
moved it up. The CEO then gave five times as much to a new hire, and twelve times as much to a third, whose answer was
right and worth 4 of the project's 22 points.

\begin{table}[!htb]
\centering
\caption{A refusal with a price quote, and what followed, from the log of one first-round live run (project p2, second
run; job j4, a level-4 logic grid worth 4 of the project's 22 points).}
\label{tab:refusal}
\footnotesize
\begin{tabular}{@{}r p{1.7cm} p{11.6cm}@{}}
\toprule
\tablehead & Given & What happened \\
\midrule
1 & \$0.001 & The CEO declares the job easy, and the hub hires \texttt{grid4} on \texttt{gpt-6-luna}. \\
2 & & \texttt{grid4}'s first call reasons up to the cap its budget allows, spends \$0.0010 and returns nothing. \\
3 & & Its next call is refused before it reaches the provider. \\
4 & & The hub moves \texttt{grid4} up to \texttt{muse-spark-1.3}, and the refusal quotes the price: ``grid4's pool has
\$0.0000, and a call on meta/muse-spark-1.3 costs about \$0.0005 to read its messages and \$0.0011 for a short answer:
reserve more.'' \\
5 & \$0.005 & The CEO checks its own budget and hires \texttt{grid4b}, which answers that the clues are inconsistent
(wrong). \\
6 & \$0.012 & The CEO hires \texttt{solve4}, whose solution is right. \\
\bottomrule
\end{tabular}
\end{table}

\textbf{The agents choose the expensive model for the right jobs, but give it too little money.} On the sixteen
projects, 65\% of the agents' tries on \texttt{muse-spark-1.3} (68 of 104) were on hard jobs, and half came after a
cheaper model had tried the same job. Its 79 hires took 55\% of what the agents spent (\$0.99 of \$1.81) and returned 12
right answers, while its measured success rate on hard jobs is 87 to 100\%. The median hire was given \$0.005, but an
attempt on the expensive model costs a median of \$0.006 to \$0.025 depending on the job type (a mean of \$0.010 to
\$0.061). So 78 of the 79 hires had a call refused because their budget could not pay for it, after spending nearly all
of it first (a median of \$0.0046). The price quote comes after the money is gone. And at $\tau = 0.5$, the budget pays
for a full attempt on the expensive model on only a few jobs, so these hires buy calls that cannot finish.

\textbf{Splitting the money pays off once the budget has room, and the cheap model alone cannot use it.} The live runs
use $\tau = 0.5$. Replaying the results table at $\tau$ of 0.25, 1 and 2 (no new model calls; Figure~\ref{fig:sweep}
and Table~\ref{tab:sweep}) shows where a strategy can win. On the sixteen projects, \texttt{gpt-6-luna} alone scores
0.740 to 0.749 at every budget. It spends 31\% of the budget at $\tau = 0.25$ and 7\% at $\tau = 2$, and has nothing to
buy with the rest. The cascade adapts during the task, as the agents do. It tries the cheapest option expected to be
good enough, moves up when a grid's check fails, and asks a third option when the two cheapest disagree. It trails the
cheap model by 0.051 at $\tau = 0.25$ and leads it by 0.133 at $\tau = 2$ ($[+0.078, +0.192]$, ahead on 12 of 16
projects). The knapsack's lead grows from 0.029 to 0.162. The best single model also changes, since the expensive model
alone passes the cheap one at $\tau = 2$ (0.815 against 0.749). The knapsack, which plans up front, leads by more than
the cascade at every budget, so the sweep does not show that adapting beats planning (V5 fails, as above). What it
shows is that a budget with room rewards a strategy that uses it. The agents ran only at $\tau = 0.5$, where the
cascade still trails the cheap model and the agents' tries on the expensive model are underfunded. We did not run them
at a larger budget, so whether they follow the cascade there is not measured, but the sweep says that a larger budget is
where they should win.

\begin{figure}[!htb]
\centering
\begin{tikzpicture}
\begin{axis}[
  width=0.78\linewidth, height=6.4cm,
  xmode=log, log basis x=2, xmin=0.2, xmax=2.5,
  xtick={0.25,0.5,1,2}, xticklabels={0.25,0.5,1,2},
  ymin=0.3, ymax=0.95, ytick={0.3,0.4,0.5,0.6,0.7,0.8,0.9}, yticklabels={0.3,0.4,0.5,0.6,0.7,0.8,0.9},
  xlabel={Budget $\tau$ (times what the best answers cost)},
  ylabel={Score (share of points)},
  label style={font=\small\color{ink}},
  tick label style={font=\footnotesize\color{muted}},
  axis line style={draw=rulecolor,line width=0.6pt}, tick style={draw=none},
  xmajorgrids, ymajorgrids, major grid style={draw=rulecolor!70,line width=0.4pt},
  every axis plot/.append style={line width=1.5pt,line join=round,line cap=round,mark size=3pt},
  legend style={font=\footnotesize,draw=none,fill=none,at={(0.5,-0.24)},anchor=north,column sep=10pt,
    cells={anchor=west}},
  legend columns=3,
]
\addplot[color=vizblue,mark=*,mark options={draw=white,line width=0.8pt,fill=vizblue}]
  coordinates {(0.25,0.696) (0.5,0.712) (1,0.783) (2,0.882)};
\addlegendentry{Cascade (adapts during the task)}
\addplot[color=vizorange,mark=square*,mark options={draw=white,line width=0.8pt,fill=vizorange}]
  coordinates {(0.25,0.747) (0.5,0.740) (1,0.748) (2,0.749)};
\addlegendentry{Every job on \texttt{gpt-6-luna}}
\addplot[color=vizaqua,mark=triangle*,mark options={draw=white,line width=0.8pt,fill=vizaqua}]
  coordinates {(0.25,0.776) (0.5,0.817) (1,0.826) (2,0.910)};
\addlegendentry{Knapsack (plans up front)}
\addplot[color=vizyellow,mark=diamond*,mark options={draw=white,line width=0.8pt,fill=vizyellow}]
  coordinates {(0.25,0.350) (0.5,0.427) (1,0.562) (2,0.815)};
\addlegendentry{Every job on \texttt{muse-spark-1.3}}
\addplot[only marks,mark=o,mark size=5pt,draw=ink,line width=1.1pt] coordinates {(0.5,0.737)};
\addlegendentry{\sys{} v2 agents (live)}
\draw[ink,line width=0.5pt] (axis cs:0.5,0.737) ++(-0.06cm,-0.19cm) -- (axis cs:0.44,0.645)
  node[anchor=north,inner sep=1pt,font=\footnotesize,text=ink,align=center,fill=white]
  {\sys{} v2 agents, live: 0.737\\(matches \texttt{gpt-6-luna})};
\end{axis}
\end{tikzpicture}
\caption{Where splitting the money pays off. Scores on the sixteen projects as the budget grows eightfold, replayed
from the results table (three runs per project). The cheap model alone stays flat, while the cascade and the knapsack
turn the extra budget into points. The agents ran live at $\tau = 0.5$ only (ringed). Table~\ref{tab:sweep} gives the
numbers and intervals.}
\label{fig:sweep}
\end{figure}

\begin{table}[!htb]
\centering
\caption{Replayed strategies on the sixteen projects at four budgets ($B = \tau\,C_\mathrm{full}$, a mean of \$0.036,
\$0.072, \$0.144 and \$0.288 a project), three runs each under v2's budget rule. The agents ran live at $\tau = 0.5$
only. The last row is the cascade's lead over the cheap model alone, with its 95\% interval over projects.}
\label{tab:sweep}
\footnotesize
\setlength{\tabcolsep}{4pt}
\begin{tabular}{@{}lrrrr@{}}
\toprule
\tablehead Strategy & $\tau = 0.25$ & $\tau = 0.5$ & $\tau = 1$ & $\tau = 2$ \\
\midrule
Every job on \texttt{gpt-6-luna} & 0.747 & 0.740 & 0.748 & 0.749 \\
Every job on \texttt{gpt-5-nano} & 0.635 & 0.700 & 0.721 & 0.723 \\
Every job on \texttt{muse-spark-1.3} & 0.350 & 0.427 & 0.562 & 0.815 \\
Equal split & 0.748 & 0.761 & 0.791 & 0.808 \\
Cascade (solver, during) & 0.696 & 0.712 & 0.783 & \textbf{0.882} \\
Knapsack (solver, up front) & 0.776 & 0.817 & 0.826 & 0.910 \\
Oracle-full (knows outcomes) & 0.825 & 0.838 & 0.844 & 0.947 \\
\sys{} v2 agents (live) & & 0.737 & & \\
\midrule
Cascade minus \texttt{gpt-6-luna} & $-0.051$ & $-0.028$ & $+0.035$ & $\mathbf{+0.133}$ \\
 & $[-0.108, +0.004]$ & $[-0.088, +0.031]$ & $[-0.036, +0.100]$ & $[+0.078, +0.192]$ \\
\bottomrule
\end{tabular}
\end{table}

\textbf{What the agents still get wrong.} Besides the expensive model, two things separate them from the knapsack.
First, the CEO keeps some jobs for itself and answers them worse than a fresh call would. In the version that hides long
sections, it answered 169 jobs itself and got 85\% right, while \texttt{gpt-6-luna} given each job on its own gets 91\%.
The 223 jobs it handed to hires were right 67\% of the time, against 68\%. Second, the difficulty the CEO declares
barely matches the jobs' hidden difficulty (rank correlation 0.10 over 128 hires), and its own calls still take 32\% of
the budget. It rarely buys from the cheap, unreliable \texttt{gridfast} (seven messages, five of them right).

Holding the other strategies to v2's rule lowers the ones that overspent under the first rule. The cascade, the hybrid
by level and the expensive model alone, which spent 1.58 to 2.14 times their budget under the first rule, fell by
0.080, 0.038 and 0.105. Random choice and the knapsack fell by only 0.002 and 0.006. It raised the cheap model alone by
0.013, because refusing a call the budget cannot pay for saves the money for later jobs (on one small project, the first
rule let the cheap model's calls spend up to 2.9 times the budget, and it scored 0.4 to 0.6 less).

\section{Limitations and Scope}
\label{sec:limitations}

\textbf{Scale of the runs.} These are first, small runs: one calibration attempt per job and option, ten projects for
the replayed allocators and six for the live ones (twice each for \sys{} v2). Replay draws on single attempts, so it
carries none of the models' own variance; the intervals we give are over projects, and with six projects they are wide.

\textbf{The calibration.} Three of the seven validity checks fail on this first table. Knowing each job's hidden
difficulty is worth under one point, so these runs cannot show what adapting during the task is worth, for any
allocator; and the dear model leaves headroom only on long aggregation jobs. A harder calibration is the first item on
the roadmap.

\textbf{Live against replay.} The baselines are replayed with each job's canonical prompt, and under the older rule that
charges a call in full even past the pool. \sys{} v2 ran live, wrote its own messages to its hires, and was held to its
pool at every call. The comparison therefore favors the replayed baselines on both counts.

\textbf{The market and the system.} Prices are list prices on the run date, and the runs have two model tiers, so local
models (and memory as a second currency) are tested in \sys{} but not measured here. Declared difficulty is the agent's
own judgment. One API key is one trust domain, and the hub's token estimate can let a run end a few percent over its
budget (2\% in one of twelve runs).

\section{Conclusion, Roadmap and Availability}
\label{sec:conclusion}

Language-model agents can now hold money and hand it to the agents they hire; until they are measured against the
best split that money could buy, how well they spend it is a matter of belief. \sys{} v2 makes agents allocate, with
reserved pools in dollars and memory, refunds on firing, and make or buy over open protocols, and records every
allocation. \bench{} makes the optimum known and holds itself to measuring allocation before any allocator runs. On
it, agents that hold the money first spent it worse than every rule we compared them with, and the traces showed
why; fixing the money mechanics closed half of the gap, and the half that remains is the agents' own judgment.

\textbf{Roadmap.} The remaining gap points at the agents rather than the mechanics. The CEO never used \texttt{attach},
so handing work over whole should not depend on the agent remembering to: a hirer could pass a job by reference by
default. The CEO spends over half its budget rereading the request on every turn; keeping long inputs out of its
context, or caching them, would leave that money for the work. Declared difficulty barely tracks real difficulty, so
the router could learn from outcomes instead of trusting the declaration. The next measurements are the ones these runs
left out: a recalibration with harder top levels and hidden bins that change success rates (checks V1, V4 and V5),
more projects and seeds, local models (and with them memory as a second currency), a small cloud tier, outside agents
whose reliability is unknown, and the cost in time of the protocol boundaries.

\textbf{Availability.} \sys{} v2 and \bench{} are source-available under the PolyForm Noncommercial License 1.0.0 at
\url{https://github.com/HASHIRU-AI/HASHIRU}, with the outcome table, every run's results and traces, and the
environment of each run. \sys{} v1's code is preserved at tag \texttt{v1.0-paper}.


\bibliographystyle{IEEEtran}
\bibliography{refs}

\end{document}
